\documentclass[sigconf,nonacm]{acmart}
% \usepackage{amsmath,amssymb,amsfonts,latexsym}
\usepackage{amsmath,amsfonts,latexsym}
\usepackage{enumerate}
\usepackage{xspace}
\usepackage{epsf,picinpar}
\usepackage{varioref}
\usepackage{colortbl,multirow,hhline}
\usepackage{listings}
\makeatletter
\let\Bbbk\relax
\makeatother
\usepackage{amssymb}
\usepackage{colortbl,multirow,hhline}
\usepackage{algorithmic}
\usepackage{algorithm}
\usepackage{caption}
\usepackage[normalem]{ulem}
\usepackage{xcolor}
\usepackage{pifont}
\usepackage{xcolor,colortbl}
\usepackage{url}
\usepackage{balance}
\usepackage{graphicx, subfigure}
\usepackage{longtable}
\usepackage{lscape}
\usepackage{multirow}
\usepackage{listings}
\usepackage{framed}
\usepackage{morefloats}
\usepackage[T1]{fontenc}
\usepackage{array}
\usepackage{pdfpages}
\usepackage{fancybox}
\usepackage{amsmath}
\usepackage{flushend}
\usepackage{booktabs}
\usepackage{enumitem}

\renewcommand{\ttdefault}{cmr}

%\newcommand{\limit}[1]{\textcolor{red}{\noindent \ding{46}~Page limit:~#1}\\}
\newcommand{\todo}[1]{\textcolor{blue}{\ding{46}~#1}} 
\newcommand{\ie}{\emph{i.e.,}\xspace}
\newcommand{\eg}{\emph{e.g.,}\xspace}
\newcommand{\etc}{etc.\xspace}
\newcommand{\etal}{\emph{et~al.}\xspace} 

\copyrightyear{2025}
\acmYear{2025}
\setcopyright{none}
%\acmConference[LD/CGT 2024/2025]{Course LD/CGT 2024/2025 - Université Libre de Bruxelles  and Vrije Universiteit Brussel}{January, 2025}{Brussels, Belgium}
\acmBooktitle{Course LD/CGT 2024/2025 - Université Libre de Bruxelles  and Vrije Universiteit Brussel, January 2025, Brussels (Belgium)}
    
\begin{document}

\title{
	{Iterative Improvement and Variable Neighborhood Descent for the Linear Ordering Problem}
}

\author{Rahman Tektas}
\affiliation{%
 \institution{542557 \\ ULB}
} \email{rahman.tektas@ulb.be}




\begin{abstract}
This report studies local search algorithms for the Linear Ordering Problem (LOP). First, iterative improvement algorithms were implemented using two pivoting rules, first-improvement and best-improvement, combined with the three neighborhoods transpose, exchange, and insert. Two initialization methods were considered: a random permutation and the Chenery–Watanabe (CW) heuristic. This yields twelve algorithmic combinations. Second, two Variable Neighborhood Descent (VND) variants were implemented, using the neighborhood orders transpose–exchange–insert and transpose–insert–exchange, with CW initialization and first-improvement only. The algorithms were evaluated on the benchmark instances provided for the course using the average percentage deviation from the best known solutions and the total computation time. Statistical comparisons were carried out with paired Student t-tests and Wilcoxon signed-rank tests in R. 
\end{abstract}

\maketitle

\section{Introduction}
The Linear Ordering Problem (LOP) is a combinatorial optimization problem defined on a square matrix C. The goal is to find a permutation of the rows and columns that maximizes the sum of matrix entries above the diagonal in the reordered matrix. Because the problem is hard, heuristic methods are often used to obtain good solutions in reasonable time.

In this work, two families of local search methods are studied. In Exercise 1.1, iterative improvement algorithms are compared under different pivoting rules, neighborhoods, and initialization methods. In Exercise 1.2, Variable Neighborhood Descent (VND) is used to combine several neighborhoods within one algorithm. The objective is to study how these design choices affect solution quality and runtime.
\input{related}
\input{definition}
\section{Results}

The benchmark instances used in the experiments were the provided LOP instances of sizes 150 and 250. Each algorithm was applied once to each instance, and a random seed was set for reproducibility and fair comparaison.
For each run, two values were recorded:
\begin{itemize}[leftmargin=1.5em]
    \item the final solution cost
    \item the computation time
\end{itemize}

Below is a summary table, where solution quality is measured using the relative percentage deviation from the best known solution, defined as
\[
\Delta_{ki} = 100 \cdot \frac{\text{best-known}_i - \text{cost}_{ki}}{\text{best-known}_i}.
\]

With this convention, smaller values correspond to better solutions.

\subsection{Exercise 1.1 : Summary of the 12 Iterative Improvement Algorithms}
\begin{table}[H]
\centering
\caption{Summary results for the 12 iterative improvement algorithms.}
\label{tab:ex1}
\begin{tabular}{lccc}
\toprule
\textbf{Algorithm Variant} & \textbf{Avg. Deviation (\%)} & \textbf{Std. Dev.} & \textbf{Total Time (s)} \\
\midrule
cw best exchange      & 3.8465  & 0.6054 & 32.9845 \\
cw best insert        & 2.2608  & 0.4065 & 33.7210 \\
cw best transpose     & 17.0783 & 1.6065 & 0.2194 \\
cw first exchange     & 2.9660  & 0.4613 & 186.7398 \\
cw first insert       & 2.0164  & 0.4109 & 128.2090 \\
cw first transpose    & 17.0863 & 1.6042 & 0.2148 \\
random best exchange  & 3.6027  & 0.5048 & 35.8138 \\
random best insert    & 2.3037  & 0.4635 & 37.9277 \\
random best transpose & 34.4701 & 3.8033 & 0.0035 \\
random first exchange & 2.8159  & 0.4810 & 228.0147 \\
random first insert   & 2.0211  & 0.4507 & 189.8870 \\
random first transpose& 34.6058 & 3.7984 & 0.0024 \\
\bottomrule
\end{tabular}
\end{table}

\subsection{Exercise 1.2 : Summary of the VND Algorithms}

\begin{table}[H]
\centering
\caption{Summary results for the two VND algorithms.}
\label{tab:ex2}
\begin{tabular}{lccc}
\toprule
\textbf{Method} & \textbf{Avg. Deviation (\%)} & \textbf{Std. Dev.} & \textbf{Total Time (s)} \\
\midrule
vnd-tei & -2.0633 & 0.4191 & 109.6544 \\
vnd-tie & -1.9959 & 0.4220 & 80.5438 \\
\bottomrule
\end{tabular}
\end{table}

\subsection{Tests}

A series of tests were performed to compare the algorithms, to answer the questions of Is there a statistically significant difference between the solution quality generated by the different
algorithms? 

The wilcoxon 




The two VND variants were compared using paired Student t-tests and Wilcoxon signed-rank tests.

\begin{itemize}[leftmargin=1.5em]
    \item Average deviation vnd-tei: \(-2.0633\)
    \item Average deviation vnd-tie: \(-1.9959\)
    \item Std deviation vnd-tei: \(0.4191\)
    \item Std deviation vnd-tie: \(0.4220\)
    \item Paired t-test p-value: \(0.2136\)
    \item Wilcoxon p-value: \(0.4169\)
\end{itemize}

Although the transpose--insert--exchange variant obtained a slightly better average deviation than the transpose--exchange--insert variant, both statistical tests returned p-values greater than \(0.05\). Therefore, no statistically significant difference in solution quality was observed between the two VND algorithms.
\input{discussion}
\section{Conclusions}\label{sec:conclusions}

In this work, iterative improvement and variable neighborhood descent algorithms were implemented for the linear ordering problem. For Exercise 1.1, twelve iterative improvement variants were evaluated by combining two pivoting rules, three neighborhoods, and two initialization methods. For Exercise 1.2, two VND variants were implemented using the required neighborhood orders and CW initialization.

The experimental results show that the choice of neighborhood has the strongest influence on solution quality. In particular, the insert neighborhood clearly outperformed the exchange neighborhood, while transpose was by far the weakest neighborhood. The results also showed that the impact of initialization depends on the strength of the neighborhood: for transpose, CW initialization was much better than random initialization, whereas for insert the difference between random and CW was not statistically significant. In addition, the insert neighborhood with first-improvement achieved better results than with best-improvement on the tested benchmark instances.

For the VND algorithms, the transpose--insert--exchange variant obtained a slightly better average deviation and a lower total runtime than the transpose--exchange--insert variant. However, the statistical tests showed that this difference in solution quality was not significant.

Overall, these experiments suggest that using a strong neighborhood such as insert is more important than using a more expensive pivoting rule, and that good initialization is especially important when the neighborhood is weak. Among the tested methods, insert-based algorithms provided the best balance between solution quality and computational effort. 

\bibliographystyle{IEEEtran}
\bibliography{references}

\end{document}

